\documentclass[a4paper,12pt]{article}
\usepackage[italian]{babel}
\usepackage[utf8]{inputenc}
\usepackage[T1]{fontenc}
\usepackage{geometry}
\usepackage{setspace}

\geometry{margin=2.5cm}
\onehalfspacing

\title{Esercizi sui Circuiti Logici}
\author{Prof. Fedeli Massimo}
\date{}

\begin{document}
	
	\maketitle
	
	Per ogni esercizio:
	
	\begin{enumerate}
		\item Scrivere la funzione logica in forma canonica (somma di mintermini)
		\item Semplificare la funzione logica
		\item Disegnare il circuito logico corrispondente
	\end{enumerate}
	
	\section*{Esercizio 1}
	
	\begin{tabular}{c c | c}
		A & B & F \\
		0 & 0 & 0 \\
		0 & 1 & 1 \\
		1 & 0 & 0 \\
		1 & 1 & 1 \\
	\end{tabular}
	
	\section*{Esercizio 2}
	
	\begin{tabular}{c c | c}
		A & B & F \\
		0 & 0 & 1 \\
		0 & 1 & 0 \\
		1 & 0 & 0 \\
		1 & 1 & 1 \\
	\end{tabular}
	
	\section*{Esercizio 3}
	
	\begin{tabular}{c c | c}
		A & B & F \\
		0 & 0 & 0 \\
		0 & 1 & 1 \\
		1 & 0 & 1 \\
		1 & 1 & 0 \\
	\end{tabular}
	
	\section*{Esercizio 4}
	
	\begin{tabular}{c c | c}
		A & B & F \\
		0 & 0 & 1 \\
		0 & 1 & 1 \\
		1 & 0 & 0 \\
		1 & 1 & 1 \\
	\end{tabular}
	
	\section*{Esercizio 5}
	
	\begin{tabular}{c c | c}
		A & B & F \\
		0 & 0 & 0 \\
		0 & 1 & 0 \\
		1 & 0 & 1 \\
		1 & 1 & 1 \\
	\end{tabular}
	
	\section*{Esercizio 6}
	
	\begin{tabular}{c c c | c}
		A & B & C & F \\
		0 & 0 & 0 & 0 \\
		0 & 0 & 1 & 1 \\
		0 & 1 & 0 & 0 \\
		0 & 1 & 1 & 1 \\
		1 & 0 & 0 & 0 \\
		1 & 0 & 1 & 1 \\
		1 & 1 & 0 & 0 \\
		1 & 1 & 1 & 1 \\
	\end{tabular}
	
	\section*{Esercizio 7}
	
	\begin{tabular}{c c c | c}
		A & B & C & F \\
		0 & 0 & 0 & 1 \\
		0 & 0 & 1 & 0 \\
		0 & 1 & 0 & 1 \\
		0 & 1 & 1 & 0 \\
		1 & 0 & 0 & 1 \\
		1 & 0 & 1 & 0 \\
		1 & 1 & 0 & 1 \\
		1 & 1 & 1 & 0 \\
	\end{tabular}
	
	\section*{Esercizio 8}
	
	\begin{tabular}{c c c | c}
		A & B & C & F \\
		0 & 0 & 0 & 0 \\
		0 & 0 & 1 & 1 \\
		0 & 1 & 0 & 1 \\
		0 & 1 & 1 & 0 \\
		1 & 0 & 0 & 1 \\
		1 & 0 & 1 & 0 \\
		1 & 1 & 0 & 0 \\
		1 & 1 & 1 & 1 \\
	\end{tabular}
	
	\section*{Esercizio 9}
	
	\begin{tabular}{c c c | c}
		A & B & C & F \\
		0 & 0 & 0 & 0 \\
		0 & 0 & 1 & 1 \\
		0 & 1 & 0 & 1 \\
		0 & 1 & 1 & 1 \\
		1 & 0 & 0 & 0 \\
		1 & 0 & 1 & 1 \\
		1 & 1 & 0 & 1 \\
		1 & 1 & 1 & 1 \\
	\end{tabular}
	
	\section*{Esercizio 10}
	
	\begin{tabular}{c c c | c}
		A & B & C & F \\
		0 & 0 & 0 & 1 \\
		0 & 0 & 1 & 0 \\
		0 & 1 & 0 & 0 \\
		0 & 1 & 1 & 1 \\
		1 & 0 & 0 & 0 \\
		1 & 0 & 1 & 1 \\
		1 & 1 & 0 & 1 \\
		1 & 1 & 1 & 0 \\
	\end{tabular}
	
	\section*{Esercizi 11--15}
	
	Per le seguenti funzioni:
	
	\begin{enumerate}
		\item scrivere la tabella di verità
		\item disegnare il circuito logico
	\end{enumerate}
	
	\begin{enumerate}
		
		\item $F(A,B,C)=A\overline{B}+BC$
		
		\item $F(A,B,C)=\overline{A}B+\overline{B}C$
		
		\item $F(A,B,C)=AB+\overline{A}C$
		
		\item $F(A,B,C)=A+B\overline{C}$
		
		\item $F(A,B,C)=\overline{A}B+AC$
		
	\end{enumerate}
	
\end{document}